\begin{center}
    {\LARGE \textbf{2007分析化学}} \\[2ex]
    {\large 时量180分钟 \quad 满分90分}
\end{center}

\vspace{0.5cm}
\noindent\textbf{注意事项:}
\begin{enumerate}[label=\arabic*., parsep=0pt, itemsep=0pt]
    \item 答题前,考生先将自己的姓名、学号、考场号、座位号填写在答题卡上,将条形码准确粘贴在考生信息条形码粘贴区。
    \item 考试结束后,将本试卷与答题纸一并收回。
    \item 回答各个问题时,将答案写在答题纸上,写在本试卷上无效。
\end{enumerate}

\vspace{0.5cm}
\noindent\textbf{一、简答题(28 pts. total)}

\begin{enumerate}[label=\arabic*., itemsep=1.5ex]
    \item 缓冲溶液的作用是什么？由一元弱酸$\ce{HA}$，其酸度常数为$K_\mathrm{a}$及其共轭碱$\ce{A^-}$组成的缓冲溶液的缓冲范围是多少？
    \item 写出下列溶液的质子条件式
    \begin{enumerate}[label=(\arabic*), noitemsep, topsep=0pt, partopsep=0pt]
        \item $0.10\ \mathrm{mol/L}\ \ce{H2SO4}$溶液
        \item 将$10\ \mathrm{mL}\ 0.10\ \mathrm{mol/L}\ \ce{NaOH}$溶液和$10\ \mathrm{mL}\ 0.20\ \mathrm{mol/L}\ \ce{HAc}$混合所得溶液
    \end{enumerate}
    \item 在进行试样分析时，常用掩蔽方法消除干扰组分的影响，常用的掩蔽方法有哪几种？
    \item 基准物质必须具备的性质是什么？EDTA可制得纯品，但EDTA标准溶液常配成大致的浓度，再进行标定，为什么？
    \item 什么是均相沉淀法，其优点是什么？
    \item 用重铬酸钾法测定铁矿石中铁的含量，加入$\ce{SnCl2}$的作用是什么？为什么要加入$\ce{H3PO4}$？
    \item 已知半反应 $\ce{VO2^+ + 2H^+ + e^- -> VO^{2+} + 2H2O}$ 的标准电极电位$E^\ominus=1.00\ \mathrm{V}$，计算$\mathrm{pH}=1.00$时溶液中电对的条件电位。
\end{enumerate}

\vspace{0.5cm}
\noindent\textbf{二、计算题(62 pts. total)}

\begin{enumerate}[label=\arabic*., itemsep=1.5ex]
    \item (12 pts.) 用$0.2\ \mathrm{mol/L}\ \ce{NaOH}$滴定$0.2\ \mathrm{mol/L}\ \ce{HCl}$和$0.2\ \mathrm{mol/L}\ \ce{NH4Cl}$的混合溶液，问能否分步滴定$\ce{HCl}$和连续滴定$\ce{NH4Cl}$？计算滴定$\ce{HCl}$的滴定误差（$\mathrm{pH_{ep}}=6.2$）。
    \item (12 pts.) 在$\mathrm{pH}=5.0$时，以二甲酚橙为指示剂用$2.00\times10^{-4}\ \mathrm{mol/L}\ \mathrm{EDTA}$滴定$2.00\times10^{-4}\ \mathrm{mol/L}$的$\ce{Pb^{2+}}$溶液，试计算调节$\mathrm{pH}$时选用六甲基四胺或$\ce{HAc-NaAc}$缓冲溶液的滴定误差各为若干？用哪种缓冲剂好？（设终点$[\ce{Ac^-}]=0.10\ \mathrm{mol/L}$，已知：$\lg K_{\mathrm{PbY}}=18.0$，$\ce{Pb(Ac)2}$的$\lg\beta_1=1.9$，$\lg\beta_2=3.3$，$K_\mathrm{a}(\ce{HAc})=1.8\times10^{-5}$；$\mathrm{pH}=5.0$时，$\lg\alpha_{\mathrm{Y(H)}}=6.6$，$\mathrm{pPb_{(二甲酚橙)}}=7.0$）
    \item (10 pts.) 采用钽试剂萃取光度法测定钒，$100\ \mathrm{mL}$试液中含钒$40\ \mathrm{mg}$，用$10\ \mathrm{mL}$钽试剂一氯仿溶液萃取，萃取率为$90\%$，用$1\ \mathrm{cm}$比色皿，在$530\ \mathrm{nm}$波长处测得萃取液的吸光度为$0.384$，求分配比和摩尔吸光系数（$A_r(\mathrm{V})=50.94$）
    \item (8 pts.) 向$20.00\ \mathrm{mL}\ 0.1000\ \mathrm{mol/L}\ \ce{Ce^{4+}}$溶液中分别加入$15.00\ \mathrm{mL}$、$20.00\ \mathrm{mL}$、$25.00\ \mathrm{mL}\ 0.1000\ \mathrm{mol/L}\ \ce{FeCl2}$，平衡时体系的电位分别为多少？（$\varphi^{\ominus\prime}_{\ce{Ce^{4+}/Ce^{3+}}}=1.44\ \mathrm{V}$，$\varphi^{\ominus\prime}_{\ce{Fe^{3+}/Fe^{2+}}}=0.68\ \mathrm{V}$）
    \item (10 pts.) 工业纯碱中总碱度的测定结果$w(\ce{Na2O})\%$为$48.45$，$48.44$，$48.42$，若标准值为$48.50\%$，计算$\bar{x}$与合格标准$48.50\%$是否有显著差异（$95\%$置信水平）？
    \\
    已知：
    \begin{center}
    \begin{tabular}{l|cc}
    $f$ & $2$ & $3$ \\
    \hline
    $t_{0.05}$ & $4.30$ & $3.18$ \\
    \end{tabular}
    \end{center}
    \item (10 pts.) 试计算$\ce{BaSO4}$在纯水中和在$\mathrm{pH}=3.0$，$\mathrm{EDTA}$浓度为$1.0\times10^{-2}\ \mathrm{mol/L}$条件下的溶解度，比较它们之间有无差异，为什么？（已知$\lg K_{\mathrm{BaY}}=7.8$；$\mathrm{pH}=3.0$时，$\lg\alpha_{\mathrm{Y(H)}}=10.8$，$\mathrm{p}K_{\mathrm{sp}}(\ce{BaSO4})=9.97$）
\end{enumerate}